% Part: computability
% Chapter: computability-theory
% Section: def-functions-self-reference

\documentclass[../../../include/open-logic-section]{subfiles}

\begin{document}

\olfileid{cmp}{thy}{slf}
\olsection{স্ব-উল্লেখ দিয়ে অপেক্ষক সংজ্ঞায়ন}

অপেক্ষককে তার নিজের সাহায্যে সংজ্ঞায়িত করতে পারা সাধারণভাবে উপযোগী।
যেমন, গণনসাধ্য অপেক্ষক $k$, $l$ ও~$m$ দেওয়া থাকলে স্থির-বিন্দু লেমা
বলে, এমন একটি আংশিক গণনসাধ্য অপেক্ষক~$f$ আছে, যা প্রত্যেক~$y$-এর জন্য
নিচের সমীকরণটি মেটায়:
\[
f(y) \simeq
\begin{cases}
k(y) & \text{যদি $l(y) = 0$} \\
f(m(y)) & \text{অন্যথায়।}
\end{cases}
\]
আরও নির্দিষ্টভাবে, $f$-কে পেতে প্রথমে
\[
g(x,y) \simeq
\begin{cases}
k(y) & \text{যদি $l(y) = 0$} \\
\cfind{x}(m(y)) & \text{অন্যথায়}
\end{cases}
\]
দিয়ে একটি অপেক্ষক সংজ্ঞায়িত করে, তারপর স্থির-বিন্দু লেমা ব্যবহার করে এমন একটি
সূচক~$e$ নেওয়া হয়, যাতে $\cfind{e}(y) \simeq g(e,y)$।
% BN-SRC-161 TARGET-CORRECTION-BEGIN
\par\smallskip\noindent\textnormal{\small\textbf{উৎস-সংশোধন (BN-SRC-161):} হিমায়িত ইংরেজি শেষ সমীকরণে সম্ভাব্য আংশিক দুই অপেক্ষকের জন্য সাধারণ “$=$” ব্যবহার করেছে। বাংলা পাঠে স্থির-বিন্দু লেমার বক্তব্য অনুযায়ী যুগপৎ সংজ্ঞায়িততার সমতা “$\simeq$” রাখা হয়েছে।} % BN-SRC-161 TARGET-CORRECTION-END

একটি নির্দিষ্ট উদাহরণ হিসেবে, “গরিষ্ঠ সাধারণ গুণনীয়ক” অপেক্ষক
$\fn{gcd}(u,v)$-কে সংজ্ঞায়িত করা যায়:
\[
\fn{gcd}(u,v) \simeq
\begin{cases}
v & \text{যদি $u = 0$} \\
\fn{gcd}(\fn{mod}(v, u), u) & \text{অন্যথায়}
\end{cases}
\]
যেখানে $\fn{mod}(v, u)$ হলো~$v$-কে~$u$ দিয়ে ভাগ করার ভাগশেষ।
স্থির-বিন্দু লেমা প্রয়োগ করে দেখা যায়, $\fn{gcd}$ আংশিক গণনসাধ্য।
(আসলে, $y$-কে জোড়া~$\tuple{u, v}$-এর সংকেত ধরে এটিকে উপরের বিন্যাসে
রাখা যায়।) এরপর~$u$-এর উপর আরোহ থেকে দেখা যায়, $\fn{gcd}$ আসলে
সর্বত্র-সংজ্ঞায়িত।

অবশ্যই, এইমাত্র আলোচিত উদাহরণগুলির চেয়ে অনেক বেশি অভিনব স্ব-উল্লেখী
সংজ্ঞা তৈরি করা যায়। অধিকাংশ প্রোগ্রামিং ভাষাই কোনো না কোনোভাবে অপেক্ষককে
তার নিজের সাহায্যে সংজ্ঞায়িত করা সমর্থন করে। লক্ষ করো, একটি অপেক্ষক
গণনাকারী অ্যালগরিদমের \emph{সূচক} ব্যবহার করে অপেক্ষকটিকে সংজ্ঞায়িত করার
সামর্থ্যের তুলনায় এটি একটু কম নাটকীয়; স্থির-বিন্দু উপপাদ্য পূর্ণ
সাধারণতায় পরের সামর্থ্যটিই দেয়।

\end{document}
